\section{Atitude}
\label{sec:cap4_atitude}

Nesta seção apresenta-se uma visão integrada da atitude do perseguidor ao longo de toda a missão de \gls{rvdb} simulada e abrange desde a aproximação de longa distância (Seção~\ref{sec:cap4_app_longa}), a aproximação de curta distância (Seção~\ref{sec:cap4_app_curta}) e a aproximação final (Seção~\ref{sec:cap4_app_final}).
O objetivo é evidenciar que o alinhamento da atitude ocorre de forma contínua desde a transferência orbital até o momento do acoplamento da ferramenta na porta de acoplamento.

Ao longo da simulação, a dinâmica de atitude do perseguidor é descrita por quaternions, a partir dos quais são obtidos os ângulos de Euler para fins de interpretação física.
Na aproximação de longa distância, a atitude é ajustada de forma a ficar igual à atitude do alvo.
Durante a aproximação de curta distância, esse alinhamento é reforçado, pois durante a transição da dinâmica de corpo rígido descrito pelas equações de Euler para a dinâmica acoplada de Lagrange, ocorre um desalinhamento por conta da movimentação do \gls{mr}.
Na aproximação final, o alinhamento de atitude continua sendo monitorado e frequentemente corrigido, pois caso haja desalinhamento, o controle das articulações do \gls{mr} irá realizar novas correções e virará um \emph{looping} de erros.

% ------------------------------------------------------------------
% Tabela: resumo global da atitude por fase da aproximação
% \begin{table}[H]
% \centering
% \caption{Resumo dos principais parâmetros de atitude ao longo das fases de aproximação.}
% \label{tab:cap4_atitude_resumo_fases}
% \begin{tabularx}{\linewidth}{|X|X|X|X|}
% \hline
% \textbf{Fase} &
% \textbf{Condição inicial de atitude} &
% \textbf{Erro máximo (ângulos de Euler / quaternion)} &
% \textbf{Condição final de atitude} \\
% \hline
% Aproximação longa & & & \\
% Aproximação curta & & & \\
% Aproximação final & & & \\
% \hline
% \end{tabularx}
% \end{table}
% ------------------------------------------------------------------

% ------------------------------------------------------------------
% Figura: evolução global dos ângulos de Euler ao longo de toda a simulação
% \begin{figure}[H]
% \centering
% \includegraphics[width=0.9\linewidth]{fig_cap4_atitude_euler_global}
% \caption{Evolução dos ângulos de Euler do perseguidor ao longo de toda a simulação (aproximações longa, curta e final).}
% \label{fig:cap4_atitude_euler_global}
% \end{figure}
% ------------------------------------------------------------------

% ------------------------------------------------------------------
% Figura: evolução dos quaternions (ou do erro em quaternion) ao longo de toda a simulação
% \begin{figure}[H]
% \centering
% \includegraphics[width=0.9\linewidth]{fig_cap4_atitude_quat_global}
% \caption{Componentes do quaternion de atitude (ou norma do erro em quaternion) do perseguidor ao longo de toda a simulação.}
% \label{fig:cap4_atitude_quat_global}
% \end{figure}
% ------------------------------------------------------------------
